\section{How to attack this paper}
\label{sec:attack}

Four labels constrain the argument. \textbf{MEASURED} means a named run or recorded observation,
with a source revision and its scope. Here it always means \emph{reported in the cited record}:
this writing task used no GPU, VM or rental. \textbf{CODE READ} means a statement checked in a
source file or a repository document; reading a reported result does not reproduce it.
\textbf{INFERRED} means an explanation or consequence, not an observation.
\textbf{UNVERIFIED} marks a question or unsupported extension; it licenses no factual assertion.
An experimental proposal below is a proposed test, not a claim that its outcome occurred.
These distinctions follow the repository's contributor rules.\src{rules}

A source key in square brackets links to an immutable repository path in the evidence index.
Unless a different revision is printed there, it means the baseline on the cover. A paragraph's
key scopes that paragraph; it does not turn a neighbouring inference into a measurement.
Where the primary host trace is unavailable, the text says the record is a secondary report.
Numbers in such a report retain that limitation even when they are precise.

\readcode The claims ratchet is \code{python3 scripts/ci/debt.py claims}. Contributors require
\code{new=0} before publishing changed documentation and may not enrol new debt to make it pass.
It is an attribution check, not a verifier of a causal explanation. In particular, its scanner
covers Markdown and Rust, rather than this paper's TeX prose. The source index and editorial
review therefore carry the TeX attribution obligation. A green ratchet alone does not certify
this paper.\src{debt}\src{rules}

\readcode The predecessor records corrections to its own claims, including a runtime-init
probe mistaken for a successful language-model workload, an object-graph hypothesis refuted by
an already available passing control, and a mutation configuration that the tool never read.
These are historical process examples, not assertions that those implementation defects remain
in v3.\src{old}

\textbf{A useful objection names a discriminating observation.} For a claim about a wake, ask
which waiter registered, which holder released, and whether the observer could see the event.
For a claim about memory, ask who owns the backing, when the host acknowledged the change and
what a failed acknowledgement leaves pending. For a speed claim, ask which revision, workload,
placement, sample unit and control produced it. A count of output tokens cannot establish text
correctness; an idle fault counter cannot prove that the workload reached the GPU.
These are methodological requirements motivated by the predecessor's correction record.\src{old}

\section{The problem and the product}

\readcode Kayfabe presents an emulated NVIDIA PCI device to a guest running NVIDIA's unmodified
driver. The host's NVIDIA driver stays loaded. The VMM process acts as a host driver client,
using authored unprivileged operations to make the real GPU perform work. The intended benefit
is simultaneous host use and guest compatibility without installing a custom forwarding driver
inside the guest. This is a research implementation with no stable installation interface or
public binary release.\src{readme}\src{architecture}\src{faq}

\readcode ``Unprivileged'' describes the host interface the design uses. It is not an
end-to-end rootless certification: the public guides say the launch scripts run as root and a
complete rootless run is not documented. ``Stock'' means no driver or kernel binary patch.
Windows still requires documented registry configuration; an unmodified binary is not a
zero-configuration guest.\src{security}\src{faq}\src{windows}

\subsection{Comparisons at the width the repository supports}

\readcode VFIO passthrough gives one guest the physical device; the host gives it up. Kayfabe
instead relies on host-authored mappings and the GPU MMU. NVIDIA vGPU is a different vendor
stack and licensing arrangement. The repository's positioning distinguishes API forwarding,
including the sister project nvkvm-pv, which uses a custom guest module, from emulating the
device for a stock guest driver. It offers no detailed Venus, virgl or native-context
performance comparison, and neither does this paper.\src{faq}

\readcode The public credits describe Caraxes, an MIT fork of Leandro, as a synthetic NVIDIA
GPU over vfio-user that plays the firmware role, shadows guest page tables and uses host
channels while keeping the host NVIDIA driver. They describe Open-IOV's goal and host-module
approach, but explicitly do not establish whether its host keeps using the GPU. These are
repository-attributed descriptions, not an independent audit of either project. The credits identify
  nvproxy as a source of capability
allowlists, not as a claim that it implements this stock-device architecture. No claim about
relative speed, security or unpublished implementation details follows.\src{credits}

\subsection{What the product argument does not establish}

\readcode A hostile-guest model is a design target, not an independent security audit. Multiple
VMs on one GPU, shared-fault containment and resource fairness need evidence of their own.
A passing Linux thin guest is not a passing Windows desktop; a passing die does not stand for
an entire family; a fast isolated renderer does not establish multi-application stability.
The public FAQ explicitly separates these cases.\src{security}\src{faq}

\inferred The compatibility surface is a moving dependency on the guest driver's observable
expectations. This is valuable when it works and costly when a seemingly unimportant ordering
or missing response sends the driver down a different path. The seven compatibility axes and
the failure case below are ways to make that cost inspectable, not a promise to eliminate it.
